\providecommand{\Obs}{\mathrm{Obs}}
\providecommand{\hCS}{\mathrm{hCS}}
\providecommand{\CE}{\mathrm{CE}}
\providecommand{\Ch}{\mathrm{Ch}}
\providecommand{\Dolb}{\mathrm{Dolb}}
\providecommand{\Conf}{\mathrm{Conf}}
\providecommand{\chir}{\mathrm{chir}}
\providecommand{\dbar}{\bar\partial}
\providecommand{\kch}{\kappa_{\mathrm{ch}}}
\providecommand{\Ran}{\mathrm{Ran}}
\providecommand{\Sym}{\mathrm{Sym}}
\providecommand{\Hom}{\mathrm{Hom}}
\providecommand{\C}{\mathbb{C}}
\providecommand{\cO}{\mathcal{O}}
\providecommand{\cE}{\mathcal{E}}
\providecommand{\cZ}{\mathcal{Z}}
\providecommand{\g}{\mathfrak{g}}
\providecommand{\HHch}{\mathrm{HH}^{\mathrm{ch}}}
\providecommand{\Rep}{\mathrm{Rep}}
\providecommand{\Alg}{\mathrm{Alg}}
\providecommand{\Fact}{\mathrm{Fact}}
\providecommand{\Hol}{\mathrm{Hol}}
\providecommand{\QME}{\mathrm{QME}}
\providecommand{\PhiFA}{\Phi^{\mathrm{FA}}}
\providecommand{\HH}{\mathrm{HH}}

\section*{Dolbeault CE cochains and local $E_3$ hCS observables on $\C^3$}

\paragraph{Local comparison.}
Let $\g$ be a finite-dimensional Lie algebra with invariant pairing and let
\[
  \cE_{\hCS}(U):=\Omega^{0,\bullet}(U,\g)
\]
for a holomorphic polydisc $U\subset \C^3$.  Classical holomorphic
Chern--Simons theory has completed local observables
\[
  \Obs^{\mathrm{cl}}_{\hCS}(U)
  =
  \widehat{\Sym}\bigl(\cE_{\hCS,c}(U)[1]^\vee\bigr),
  \qquad
  d_{\mathrm{BV}}=\dbar+d_{\CE}.
\]
The Dolbeault--chiral Chevalley--Eilenberg model is
\[
  \CE^*_{\dbar,\chir}\!\bigl(\cE_{\hCS}|_U,\cO_U\bigr)
  =
  \Bigl(
    \prod_{n\geq 0}
    \Omega^{0,\bullet}\!\bigl(U^n,
      \Hom(\Sym^n_{\chir}\cE_{\hCS},\Delta^{(n)}_*\cO_U)\bigr)^{\Sigma_n},
    \dbar+d_{\chir}
  \Bigr).
\]
For the current Lie algebroid $\Omega^{0,\bullet}(-,\g)$ the chiral
bracket is the diagonal distribution
\[
  \mu(\alpha,\beta)=[\alpha,\beta]_\g\,\delta_\Delta ,
\]
and $[\dbar,d_{\chir}]=0$ because $\dbar$ acts on Dolbeault form degree
while $[-,-]_\g$ acts on coefficients.  Hence
\[
  \Obs^{\mathrm{cl}}_{\hCS}(U)
  \simeq
  \CE^*_{\dbar,\chir}\!\bigl(\cE_{\hCS}|_U,\cO_U\bigr)
\]
as complete filtered cochain complexes, compatibly with disjoint-open
factorisation maps.  This is a classical comparison; the quantum complex
is obtained only after a Costello renormalisation datum solving the
quantum master equation.

\paragraph{Holomorphic $E_3$ structure.}
The assignment
\[
  U\longmapsto \Obs_{\hCS}(U)
\]
is a holomorphic factorisation algebra on $\C^3$.  Gwilliam--Williams
identify holomorphic factorisation algebras on $\C^n$ with algebras over
the holomorphic little-$n$-polydisc operad in $\Ch(\Dolb)$.  For $n=3$
the local disk value therefore carries a holomorphic $E_3$ action:
\[
  \Obs_{\hCS}(D^3)\in
  \Alg_{E^{\Hol}_3}\bigl(\Ch(\Dolb)\bigr).
\]
This $E_3$ structure is native to the local factorisation algebra of
observables.  It is distinct from the CY-to-chiral specialised output
$A=\Phi_3^{(\Sigma_2,C)}(\cC)$, which is $E_1$-chiral in Vol~III.  The
$E_2$-braided object belongs to the derived or Drinfeld centre
\[
  \cZ\bigl(\Rep^{E_1}(A)\bigr)
\]
or to a Drinfeld double after the Hopf pairing and completion have been
constructed.

\paragraph{The deformation-complex shift.}
Kontsevich, Tamarkin, and Francis prove that the deformation complex of an
$E_n$-algebra is naturally $E_{n+1}$.  Applied to an $E_3$ algebra $B$,
this produces an $E_4$ algebra
\[
  \mathrm{Def}_{E_3}(B)\simeq \HH_{E_3}(B,B)[3].
\]
For $B=\cO_{\C^3}$ this object controls deformations of the commutative
base algebra.  It is not the hCS observable algebra above.  The
target comparison is
\[
  \Obs_{\hCS}
  \simeq
  \CE^*_{\dbar,\chir}(\cE_{\hCS},\cO_{\C^3})
\]
inside holomorphic factorisation algebras; the $E_4$ deformation complex
is a separate tangent object.

\paragraph{Binary product on CE cochains.}
For disjoint polydiscs $D_1,D_2\subset D\subset\C^3$ and
\[
 f_i\in
 \Hom_{\cD_{D_i^{n_i}}}
 \bigl(\Sym^{n_i}_{\chir}\cE_{\hCS}|_{D_i},\Delta_*\cO_{D_i}\bigr)
\]
of Dolbeault type $(0,q_i)$, the factorisation product is
\[
  \mu_{D_1,D_2}(f_1\otimes f_2)(l_1,\ldots,l_{n_1+n_2})
  =
  \sum_{\sigma\in \operatorname{Sh}(n_1,n_2)}
  \epsilon(\sigma)\,
  f_1(l_{\sigma(1)},\ldots,l_{\sigma(n_1)})\,
  f_2(l_{\sigma(n_1+1)},\ldots,l_{\sigma(n_1+n_2)}),
\]
after restriction of the first block to $\Conf_{n_1}(D_1)$ and the
second block to $\Conf_{n_2}(D_2)$.  The product has Dolbeault type
$(0,q_1+q_2)$.  Coherent higher products are parametrised by chains on
$\Conf_k(\C^3)$; collision strata are governed by the chiral bracket and,
in the quantum theory, by the renormalised BV propagator.

\paragraph{Quantum correction.}
In the abelian branch $\g=\C$ the bracket vanishes and the quantum
differential is the free Costello differential
\[
  D_\hbar=\dbar+\hbar\Delta_{\mathrm{ren}},
\]
with no local interaction counterterm.  For non-abelian $\g$ the statement
requires the full QME package:
\[
  D_\hbar
  =
  \dbar+d_{\chir}
  +\hbar\Delta_{\mathrm{ren}}
  +\sum_{r\geq 1}\hbar^r I_r,
  \qquad
  D_\hbar^2=0.
\]
The terms $I_r$ are local Costello counterterms.  Their existence is an
input to the quantum comparison, not a consequence of the classical CE
identification.

\paragraph{Koszul layer.}
After augmentation, conilpotence, and completion hypotheses, the local
CE algebra and the local $E_3$ enveloping algebra are Koszul dual:
\[
  U^{\Hol}_{E_3}(\cE_{\hCS})
  \simeq
  \Omega_{E_3}
  \CE^*_{\dbar,\chir}(\cE_{\hCS},\C),
  \qquad
  \CE^*_{\dbar,\chir}(\cE_{\hCS},\C)
  \simeq
  B_{E_3}U^{\Hol}_{E_3}(\cE_{\hCS}).
\]
This is the $A^!$ layer.  It is separate from the algebra $A$, the bar
object $B(A)$, the cooperadic dual $A^i$, and the derived chiral centre
$Z^{\mathrm{der}}_{\mathrm{ch}}(A)$.

\paragraph{Consequences for $\Phi$ and $\kch$.}
For compact CY$_3$ targets the canonical Vol~III scalar attached to the
constructed $\PhiFA_3$ locus is
\[
  \kch(A_X)=\sum_q(-1)^q h^{0,q}(X).
\]
The CE Euler characteristic of a gauge-theory model computes this same
number only after the hCS model has been identified with the normalised
$\PhiFA_3$ output and the coefficient sector has superdimension one.  In
particular, no extra factor such as
$\dim \CE^*(\g)_{\mathrm{hor}}$ belongs to the canonical
$\kch(A_X)$ formula.  For the abelian local model on $\C^3$ the
normalised value is $\kch(\C^3)=1$.

\paragraph{Proof obligations.}
The local classical hCS/CE comparison on $\C^3$ is chain-level.  The
non-abelian quantum comparison requires a chosen gauge fixing, heat-kernel
renormalisation, QME anomaly cancellation, completed observables, and
compatibility of counterterms with factorisation maps.  Compact CY$_3$
closure further requires a TCFT or framing datum, anomaly control, the
completion of global observables, and the relevant negative-Hochschild
filtration.

\paragraph{References.}
\begin{itemize}
  \item M.\ Kontsevich, ``Operads and motives in deformation
    quantization,'' \emph{Lett.\ Math.\ Phys.} 48 (1999), 35--72, \S4.
  \item D.\ Tamarkin, ``Deformation complex of a $d$-algebra is a
    $(d{+}1)$-algebra,'' arXiv:math/0010072.
  \item J.\ Francis, ``The tangent complex and Hochschild cohomology
    of $E_n$-rings,'' \emph{Compos.\ Math.} 149 (2013), 430--480.
  \item J.\ Francis and D.\ Gaitsgory, ``Chiral Koszul duality,''
    \emph{Selecta Math.} 18 (2012), 27--87.
  \item O.\ Gwilliam and B.\ R.\ Williams, ``Higher Kac--Moody
    algebras and symmetries of holomorphic field theories,''
    arXiv:2009.05037.
  \item K.\ Costello and O.\ Gwilliam, \emph{Factorization Algebras
    in QFT}, Vols.~I--II.
\end{itemize}
